\section{A compact contraction of one plane}\label{sec:isolated-plane}

The isolated example separates the two ingredients of the construction.
The Mukai square bound controls the entire compactification, while the
three-dimensional extension space identifies the exceptional plane.
Its topology and metric can then be calculated in the same normalization.

\subsection{The surface and the central factors}
Take the even lattice
\[
 \NS(X)=\Z H\oplus\Z D,\qquad H^2=2,\quad H\cdot D=0,
 \quad D^2=-22.
\]
Morrison's realization theorem \cite[Corollary~2.9(i)]{morrison}
gives a projective K3 surface with precisely this N\'eron--Severi
lattice. A root would satisfy $x^2-11y^2=-1$, which is impossible
modulo eleven. The positive cone is therefore an ample chamber;
choose its component so that $H$ is ample. Put
\begin{equation}\label{eq:isolated-vectors}
 a=(1,0,1),\qquad b=(2,D,-5),\qquad v=a+b=(3,D,-4).
\end{equation}
We have $a^2=b^2=-2$, $(a,b)=3$, $v^2=2$ and
$c_2(b)=c_2(v)=-4$.

\begin{lemma}\label{lem:isolated-factors}
There is a unique $H$-slope-stable bundle $B$ with vector $b$.
It is rigid. Every central $H$-polystable class of vector $v$
is either stable or $\OO_X\oplus B$. Every $H$-semistable
torsion-free sheaf of vector $v$ is locally free. No central ideal
object with that total vector has positive lost charge.
\end{lemma}

\begin{proof}
Yoshioka's primitive-vector theorem
\cite[Theorem~0.1 and condition $(\natural)$]{yoshioka99}
gives a sheaf of vector $b$ at a generic polarization arbitrarily
near $H$. Choose the polarization in a chamber adjoining $H$ for
the finite slope-wall arrangement
\cite[Theorem~6.6]{grebtoma}. Generic slope equality would require
$D/2$ to be integral. It is not, so the sheaf is slope stable.
As the polarization tends to $H$ within that chamber, every fixed
subsheaf inequality becomes weak; the sheaf is $H$-semistable.

An $H$-slope-zero rank-one torsion-free factor has the form
$\OO_X(kD)\otimes I_T$, with vector
$(1,kD,1-11k^2-n)$, where $n=\operatorname{length}T\ge0$.
If a sheaf of vector $b$ split into two such factors, its second
Chern class would be
\[
 22k(k-1)+n_1+n_2\ge0,
\]
contrary to $c_2(b)=-4$. Thus it is $H$-stable.

For any slope-stable torsion-free sheaf $F$, its double dual is
stable: intersect a subsheaf of $F^{**}$ with $F$ to preserve
rank and first Chern class. If $l=\operatorname{length}(F^{**}/F)$,
then
\begin{equation}\label{eq:double-dual-square}
 v(F^{**})^2=v(F)^2-2\rank(F)l.
\end{equation}
A simple sheaf on a K3 surface has Mukai square at least $-2$,
by Riemann--Roch and Serre duality. For $b^2=-2$ this forces
$l=0$. Two nonisomorphic stable bundles with vector $b$ would
have $\Hom=\Ext^2=0$, whereas their Euler pairing is two. This
is impossible. The same calculation for $B$ itself gives
$\Ext^1(B,B)=0$.

Now allow an ideal object with residual vector
$v_l=(3,D,-4+l)$, where $l\ge0$ is its total lost charge.
The residual stable factors have slope zero. A stable rank-three
factor has square $2-6l\ge-2$, hence $l=0$. For a rank partition
$1+2$, take the rank-one vector
$c=(1,kD,1-11k^2-n)$. The complementary stable rank-two vector
has square
\begin{equation}\label{eq:isolated-rank-two-bound}
 (v_l-c)^2=-2+44k-66k^2-4(n+l).
\end{equation}
For integral $k$, the quadratic term is nonpositive, with equality
only at zero. The lower square bound forces $k=n=l=0$. The
factors are $\OO_X$ and $B$.

For three rank-one factors, write $\sum k_i=1$. Their second
Chern class is
\[
 11\left(\sum_i k_i^2-1\right)+\sum_i n_i\ge0,
\]
whereas $c_2(v_l)=-4-l<0$. These are all positive-rank
partitions. Finally a stable sheaf of vector $v$ is locally free
by \eqref{eq:double-dual-square}, and a strictly semistable one
is an extension of the two locally free factors. Such an
extension is locally free.
\end{proof}

We will also need the following consequence of surface gauge
compactness, which applies to both examples in this paper.

\begin{lemma}\label{lem:no-charge-compactness}
Fix a topological bundle type on a K3 surface and a smooth compact
family of base K\"ahler metrics through $\omega_0$. Keep the
holomorphic determinant fixed and let its HE connection vary
smoothly. Suppose every central ideal HYM object of that total
type has zero lost charge. A sequence of HYM connections whose
parameters tend to zero then has a subsequence converging
smoothly modulo gauge on all of $X$. Moreover, after shrinking
the parameter neighborhood, no ideal object of that total type
has positive lost charge at any parameter in it.
\end{lemma}

\begin{proof}
The trace-free connections are projective ASD connections, and
their energies are topologically bounded. Pass to a weak limit
of their curvature measures and remove the finitely many atoms
above the small-energy threshold. In fixed coordinate balls
away from those points, the base metrics converge smoothly.
Their Sobolev, gauge-fixing and elliptic constants are uniformly
bounded. Small-energy regularity first gives local curvature
bounds stronger than $L^2$; varying-metric compactness and
elliptic bootstrapping then give smooth subsequential convergence
in gauges on the complement. This is the setting of
\cite[Introduction, Theorem~E$'$, and Theorem~10.3]{wehrheim}.
An energy bound alone would not meet that theorem's stronger
local integrability hypothesis.

The surface compactness, extension, removal and defect-current
results \cite[Section~3.1, Theorems~3.3, 3.5--3.7 and~3.10]{gstw}
produce a central polystable residual bundle and a nonnegative
integral lost charge. The determinant is smooth and has no
concentrating part. Local projective instantons lift over
four-balls, so the missing charges are integers. By hypothesis
their sum is zero. The energy identity and small-energy
regularity therefore remove all the exceptional points, proving
smooth convergence on $X$.

If positive-charge ideal objects occurred arbitrarily close to
zero, choose such a sequence. The energy bound permits a
subsequence with one fixed positive defect. Applying compactness
to its residual connections can add only nonnegative lost
charge. The resulting central ideal object would still have
positive defect, a contradiction. This also accounts for the
zero-dimensional discrepancy of a torsion-free sheaf from its
polystable double dual.
\end{proof}

\subsection{The two chambers and their full exceptional planes}
Let $\omega_t$ be the Ricci-flat representative of $H+tD$.
The base volume is $1-11t^2$. Shrink the interval so that no
other slope wall for $b$ or $v$ meets it. For large integral
$m$, the degree of $v$ at $mH\pm D$ is $\mp22$, coprime to
three. Thus its Gieseker semistable sheaves are slope stable.
Yoshioka's theorem and its deformation-type refinement
\cite[Appendix, Theorem~8.1]{yoshioka01} give smooth projective
fourfolds
\[
 M_\pm=M_{mH\pm D}(v)
\]
of K3$^{[2]}$ type. All their points are bundles, by
\eqref{eq:double-dual-square}. Since $\chi(E)=-1$, the
universal-family criterion \cite[Theorem~4.6.5]{huybrechtslehn}
gives ordinary universal families.

Every chamber object is $H$-semistable, by taking limits of its
subsheaf inequalities; an $H$-stable object remains stable in
both adjacent chambers. The only remaining central graded type
is $\OO_X\oplus B$. Cross Hom and second Ext groups vanish,
and both cross first Ext groups have dimension three. At $t>0$
the allowable orientation is
\begin{equation}\label{eq:isolated-positive-extension}
 0\longrightarrow B\longrightarrow E\longrightarrow\OO_X
 \longrightarrow0,
\end{equation}
since $\mu_t(B)=-11t<\mu_t(E)=-22t/3$. The orientation is
reversed at $t<0$.

The full local model at $E_0=\OO_X\oplus B$ has coordinates
$(u,\varphi)\in\C^3\oplus(\C^3)^*$, with action
$(u,\varphi)\mapsto(e^{i\theta}u,e^{-i\theta}\varphi)$,
complex equation $\varphi(u)=0$, and real level
\begin{equation}\label{eq:isolated-level}
 |u|^2-|\varphi|^2=\tau(t),\qquad
 \tau(t)=-\frac{44\pi}{3}t.
\end{equation}
Here $u$ is in $\Ext^1(B,\OO_X)$. Formula
\eqref{eq:general-real-level} gives the sign and coefficient.
The determinant-one infinitesimal lift has weights $2/3,-1/3$,
whose difference is one; there is no additional rank factor in
the effective circle moment.

For $\tau>0$ the complex quotient is $T^*\PP^2$ with
coordinate $[u]$ on the zero section; for $\tau<0$ use
$[\varphi]$. The affine map is
\begin{equation}\label{eq:isolated-incidence}
 (u,\varphi)\longmapsto u\otimes\varphi,
 \qquad \{A\in\mathfrak{sl}_3:A^2=0,\ \rank A\le1\}.
\end{equation}
The entire zero section has coefficient norm squared $|\tau|$.
It therefore lies in the common coefficient neighborhood for
small $t$. Corollary~\ref{cor:harmonic-comparison} gives actual
stable bundles on that entire section. The coordinate change
fixes each one-direction extension axis: its equivariance and
identity derivative exclude higher monomials of weight one on
that axis. The harmonic recursion is also linear there, because
the square of a one-direction extension is zero. Thus these are
the actual projective extension families in
\eqref{eq:isolated-positive-extension} and its reverse.

For clarity, the same chart contains every point of the actual
central fibre. If not, take $t_n\to0$ and bundles outside it
with graded class $\OO_X\oplus B$. Lemmas
\ref{lem:isolated-factors} and \ref{lem:no-charge-compactness}
give a smooth gauge limit $G$. Normalize the inclusion of the
fixed first factor in each extension to have $L^2$ norm one.
Its holomorphic-map equation has smoothly convergent coefficients;
global elliptic estimates give smooth convergence to a nonzero
map into $G$. Equal-slope stability and polystability split off
that factor. The remaining factor is identified by its vector
and the uniqueness just proved. Thus $G=\OO_X\oplus B$, which
puts the sequence in the open actual solution chart, a
contradiction. The extension planes are embedded: the
classifying map is biholomorphic on the smooth quotient, and
the middle term determines its distinguished subobject up to
scalar. Write their images as $P_\pm$.

\subsection{The common normal target}
The smooth chamber bases are weakly normal, so the
Greb--Toma classifying theorem applies to their universal
families \cite[Main Theorem]{grebtoma}. For a surface the
slope equivalence is given by the graded double dual and its
defect cycle \cite[Theorem~5.5, Proposition~5.8 and
Remark~5.9]{grebtoma}. Our maps have a common reduced projective
image: its points are the common stable open set and the one
class $z_0=\OO_X\oplus B$. Their fibres are points and the
entire planes just identified.

Normalize this reduced image. Stein factorization of either
proper birational map has a normal finite factor. Its finite
fibres have one point because the original fibres are connected.
It is therefore the normalization, bijective over the reduced
image. The two chamber maps lift to a common normal projective
target $Z$, with the same point topology and the same fibres.

We explain the analytic identification, since a topological
description of the fibres would not suffice. The local
holomorphic classifying family commutes with central
S-equivalence, which in the finite model is exactly the
invariant $u\otimes\varphi$. Its chart contains the full
exceptional plane. Properness now allows the target to be
shrunk until its entire inverse image is in that chart:
otherwise a sequence outside the chart mapping to $z_0$
would have a limit in the plane. The proper incidence map
in \eqref{eq:isolated-incidence} has the same property.

Both targets are normal. For the finite target, the hypersurface
$\varphi(u)=0$ is an irreducible quadric with an isolated
singularity, hence normal, and its invariant ring is normal.
A proper birational map to a normal target has direct image
of its structure sheaf equal to that structure sheaf. Pull
back target coordinate functions to the two identified
resolution neighborhoods and descend them by this identity.
The resulting maps between target germs are inverse on the
dense stable open set, hence everywhere. This identifies the
actual contraction with \eqref{eq:isolated-incidence}.

Consequently $z_0$ is the unique singular point and both maps
are small symplectic resolutions. Zero-level points with
$u,\varphi\ne0$ and $\varphi(u)=0$ are stable and tend to
the origin. The listed target is therefore the closure of
the stable wall locus. The surface comparison
\cite[Theorem~4.11 and Corollary~4.15]{gstw} identifies its
compact topology with the HYM compactification; its analytic
germ was proved separately above.

At the central polarization, the ordinary Gieseker moduli
is a different object. Its reduced Euler characteristic is
$-1/3$, which cannot equal that of a factor of rank one or
two. The nonsplit extensions in
\eqref{eq:isolated-positive-extension} are Gieseker stable
at $H$, since the $B$ subobject has reduced constant term
$-3/2$; the opposite orientation is destabilized by
$\OO_X$. Thus $M_H(v)$ is the resolution $M_+$, not $Z$.
The distinction between the slope target and the central
Gieseker space is already visible in this simplest example.

The construction has now placed the full plane, with its
cotangent normal bundle and its normalized quotient metric,
inside a compact fourfold. Sections~\ref{sec:integral-attachments}
and \ref{sec:harmonic-limits} compute its integral attachment
and metric collapse. They will use the same extension family
and the same signed Mukai marking.

